\section{Procesos de Soporte}

Esta sección describe los procesos transversales que soportan la ejecución del proyecto Kit Pagos Colombia, para que todos los integrantes del equipo conozcan cuáles son estos procesos y cuáles son sus responsabilidades específicas en cada uno de ellos. Estos procesos abarcan desde la regulación del trabajo en equipo y la administración de riesgos, hasta el control de cambios, la medición del progreso y las actividades de aseguramiento de calidad.

Los procesos de soporte se organizan en cinco áreas funcionales interrelacionadas: ambiente de trabajo (reglas de convivencia y mecanismos de control), administración de riesgos (identificación, valoración y mitigación), administración de configuración y documentación (control de cambios e ítems de configuración), métricas y medición (recolección y consolidación de datos de desempeño), y control de calidad (verificación, validación y auditorías). Cada área se ejecuta de manera continua a lo largo del proyecto, integrándose en los ciclos de iteración DSR definidos en la sección 3.1.

\subsection{Ambiente de Trabajo}

El equipo Kit Pagos Colombia trabaja bajo un conjunto explícito de reglas diseñadas para asegurar que la comunicación sea clara, que el trabajo esté siempre documentado y que las responsabilidades sean bien distribuidas. El equipo se reúne con el director los lunes de 2 a 3 p.~m., por Microsoft Teams, y se coordina a diario por un grupo de WhatsApp. La versión 1.0 de este plan preveía dos reuniones internas por semana (los lunes a las 10 a.~m. y los jueves a las 3 p.~m.), actas en Notion y un canal de Discord. Según el equipo, nada de eso se usó. Las decisiones y los compromisos quedan en los issues, en los pull requests y en el registro de arquitectura. Se espera que todos los integrantes respondan consultas dentro de máximo 24 horas hábiles.

Las tareas se asignan como issues en el hito de la semana: cada iteración de dos semanas tiene un hito por semana en GitHub, y cada issue nombra a su responsable. Ningún integrante puede asumir simultáneamente más de 40 horas de trabajo semanal en el proyecto. Todo código debe compilar en TypeScript estricto y cumplir las reglas de ESLint del paquete (\texttt{eslint.config.mjs}); el plan original remitía a unos estándares del «documento técnico del SDK (sección 1.4)», pero la sección 1.4 de este plan son los entregables. Antes de integrarse a \texttt{devops}, cada pull request pasa la integración continua, que exige la cobertura mínima del 80\% y las métricas CK. Además, lo aprueba un integrante distinto al autor (condición 1 de la definición de terminado, \texttt{methodology.md}, sección 6). Los documentos también se revisan en pull requests, con comentarios en GitHub. Los artefactos que no cumplan con estos criterios de calidad son devueltos para corrección sin excepción.

La puntualidad en las entregas es crítica. Las entregas de artefactos se realizan en las fechas comprometidas al inicio de cada ciclo. Si un integrante prevé que no podrá cumplir una entrega dentro de 48 horas, notifica de inmediato al Coordinador del equipo y al director del proyecto, explicando la causa y proponiendo una nueva fecha de entrega. Las decisiones sobre extensión de plazo las toma el director del proyecto. En caso de desacuerdo técnico entre dos o más integrantes sobre arquitectura, diseño o decisiones de implementación, cada uno deja su posición en el issue o en el pull request. El Líder de arquitectura y núcleo SDK o el Coordinador arbitran para tomar una decisión consensuada. Si la discrepancia persiste, se escala al director del proyecto. Se preserva en todo momento el respeto y la profesionalidad; los conflictos técnicos nunca deben afectar las relaciones personales del equipo.

La versión 1.0 preveía también un reporte de completitud cada viernes, con una reunión extraordinaria si bajaba del 60\%, y un sumario de decisiones por fase en Notion. No se hicieron. El avance se lee en los hitos de GitHub. Las decisiones arquitectónicas, los cambios de alcance y lo medido quedan en el registro de arquitectura con su fecha y su responsable, lo que permite auditar la trazabilidad histórica de las decisiones.
